% GENERATED FILE. Edit canonical lessons, then run npm run generate:books.
% canonical-source-hash: fnv1a64:230206286ecbe2d8
% canonical-lessons: ES-C657-caducar, ES-C657-valido, ES-C657-renovar, ES-C657-fotocopia, ES-C657-consulado

\chapter{Papeles en regla}
\label{ch:es-papeles-en-regla}
\hlchaptermodality{657}

\begin{chapteropening}
\textbf{By the end of this chapter:} \emph{I can check that my documents are valid, renew a passport, bring a photocopy and go to the consulate.}

Complete the last lesson of chapter 657: I can check that my documents are valid, renew a passport, bring a photocopy and go to the consulate.
\end{chapteropening}

\section[caducar]{caducar — to expire}

\label{lesson:ES-C657-caducar}



\begin{quote}
Before the new one: say the Spanish for to stamp, then the Spanish for to search.
\end{quote}

\subsection*{You'll want to know: caducar}
\textbf{caducar} — \textquotedblleft{}to expire\textquotedblright{}.

\begin{quote}
Mi pasaporte caduca en marzo.
\end{quote}

\begin{quote}
\emph{My passport expires in March.}
\end{quote}

Food \emph{caduca} too: \emph{la fecha de caducidad} is the use-by date.

\begin{cousinweb}
\emph{caducar} is built on \emph{caduco} \textquotedblleft{}worn out, past its time\textquotedblright{} (EtymDB), a learned borrowing from Latin \emph{cadūcus} \textquotedblleft{}falling, fallen\textquotedblright{} (Lewis \& Short), from \emph{cadere} \textquotedblleft{}to fall\textquotedblright{}.
\end{cousinweb}

\begin{grammarlens}[title={What you've built}]
One more word for telling, explaining and describing at B1.
\end{grammarlens}

\subsection*{Your turn}
\begin{itemize}
\raggedright
  \item \emph{Say it:} \emph{caducar}
  \item \emph{Say it:} \emph{caducar}, once more
  \item \emph{Say it:} a sentence of your own with \emph{caducar}
  \item \emph{Recall:} say the Spanish for to stamp, then the Spanish for to search, then say \emph{caducar} again
\end{itemize}

\begin{tcolorbox}[breakable,colback=teal!4,colframe=teal!35!black,title={Before you move on}]
What does caducar mean? (\textquotedblleft{}To expire\textquotedblright{}.) Say it once more.
\end{tcolorbox}
\section[válido]{válido — valid}

\label{lesson:ES-C657-valido}



\begin{quote}
Before the new one: say the Spanish for to search, then the Spanish for to expire.
\end{quote}

\subsection*{You'll want to know: válido}
\textbf{válido} — \textquotedblleft{}valid\textquotedblright{}.

\begin{quote}
Este billete es válido por un mes.
\end{quote}

\begin{quote}
\emph{This ticket is valid for a month.}
\end{quote}

\emph{Válido hasta} — valid until.

\begin{cousinweb}
\emph{válido} is a learned borrowing from Latin \emph{validus} \textquotedblleft{}strong, able, powerful\textquotedblright{} (Lewis \& Short; EtymDB), from \emph{valēre} \textquotedblleft{}to be strong\textquotedblright{}.
\end{cousinweb}

\begin{grammarlens}[title={What you've built}]
One more word for telling, explaining and describing at B1.
\end{grammarlens}

\subsection*{Your turn}
\begin{itemize}
\raggedright
  \item \emph{Say it:} \emph{válido}
  \item \emph{Say it:} \emph{válido}, once more
  \item \emph{Say it:} a sentence of your own with \emph{válido}
  \item \emph{Recall:} say the Spanish for to search, then the Spanish for to expire, then say \emph{válido} again
\end{itemize}

\begin{tcolorbox}[breakable,colback=teal!4,colframe=teal!35!black,title={Before you move on}]
What does válido mean? (\textquotedblleft{}Valid\textquotedblright{}.) Say it once more.
\end{tcolorbox}
\section[renovar]{renovar — to renew}

\label{lesson:ES-C657-renovar}



\begin{quote}
Before the new one: say the Spanish for to expire, then the Spanish for valid.
\end{quote}

\subsection*{You'll want to know: renovar}
\textbf{renovar} — \textquotedblleft{}to renew\textquotedblright{}.

\begin{quote}
Tengo que renovar el pasaporte.
\end{quote}

\begin{quote}
\emph{I have to renew my passport.}
\end{quote}

The o becomes ue under stress: \emph{renuevo}, \emph{renueva}.

\begin{cousinweb}
\emph{renovar} comes from Latin \emph{renovāre} \textquotedblleft{}to renew, to restore\textquotedblright{} (Lewis \& Short), built on \emph{novus} \textquotedblleft{}new\textquotedblright{} (EtymDB).
\end{cousinweb}

\begin{grammarlens}[title={What you've built}]
One more word for telling, explaining and describing at B1.
\end{grammarlens}

\subsection*{Your turn}
\begin{itemize}
\raggedright
  \item \emph{Say it:} \emph{renovar}
  \item \emph{Say it:} \emph{renovar}, once more
  \item \emph{Say it:} a sentence of your own with \emph{renovar}
  \item \emph{Recall:} say the Spanish for to expire, then the Spanish for valid, then say \emph{renovar} again
\end{itemize}

\begin{tcolorbox}[breakable,colback=teal!4,colframe=teal!35!black,title={Before you move on}]
What does renovar mean? (\textquotedblleft{}To renew\textquotedblright{}.) Say it once more.
\end{tcolorbox}
\section[la fotocopia]{la fotocopia — the photocopy}

\label{lesson:ES-C657-fotocopia}



\begin{quote}
Before the new one: say the Spanish for valid, then the Spanish for to renew.
\end{quote}

\subsection*{You'll want to know: la fotocopia}
\textbf{la fotocopia} — \textquotedblleft{}the photocopy\textquotedblright{}.

\begin{quote}
Trae una fotocopia del pasaporte.
\end{quote}

\begin{quote}
\emph{Bring a photocopy of your passport.}
\end{quote}

The verb is \emph{fotocopiar}.

\begin{cousinweb}
\emph{fotocopia} is built from Greek \emph{phōs, phōtós} \textquotedblleft{}light\textquotedblright{} and \emph{copia} \textquotedblleft{}copy\textquotedblright{}, which comes from Latin \emph{cōpia} \textquotedblleft{}plenty\textquotedblright{}.
\end{cousinweb}

\begin{grammarlens}[title={What you've built}]
One more word for telling, explaining and describing at B1.
\end{grammarlens}

\subsection*{Your turn}
\begin{itemize}
\raggedright
  \item \emph{Say it:} \emph{la fotocopia}
  \item \emph{Say it:} \emph{la fotocopia}, once more
  \item \emph{Say it:} a sentence of your own with \emph{la fotocopia}
  \item \emph{Recall:} say the Spanish for valid, then the Spanish for to renew, then say \emph{la fotocopia} again
\end{itemize}

\begin{tcolorbox}[breakable,colback=teal!4,colframe=teal!35!black,title={Before you move on}]
What does la fotocopia mean? (\textquotedblleft{}The photocopy\textquotedblright{}.) Say it once more.
\end{tcolorbox}
\section[el consulado]{el consulado — the consulate}

\label{lesson:ES-C657-consulado}



\begin{quote}
Before the new one: say the Spanish for to renew, then the Spanish for the photocopy.
\end{quote}

\subsection*{You'll want to know: el consulado}
\textbf{el consulado} — \textquotedblleft{}the consulate\textquotedblright{}.

\begin{quote}
El consulado abre por la mañana.
\end{quote}

\begin{quote}
\emph{The consulate opens in the morning.}
\end{quote}

The person is \emph{el cónsul}.

\begin{cousinweb}
\emph{consulado} is a learned borrowing from Latin \emph{cōnsulātus} \textquotedblleft{}the office of consul\textquotedblright{} (Lewis \& Short; EtymDB), from \emph{cōnsul}.
\end{cousinweb}

\begin{grammarlens}[title={What you've built}]
One more word for telling, explaining and describing at B1.
\end{grammarlens}

\subsection*{Your turn}
\begin{itemize}
\raggedright
  \item \emph{Say it:} \emph{el consulado}
  \item \emph{Say it:} \emph{el consulado}, once more
  \item \emph{Say it:} a sentence of your own with \emph{el consulado}
  \item \emph{Recall:} say the Spanish for to renew, then the Spanish for the photocopy, then say \emph{el consulado} again
\end{itemize}

\begin{tcolorbox}[breakable,colback=teal!4,colframe=teal!35!black,title={Before you move on}]
What does el consulado mean? (\textquotedblleft{}The consulate\textquotedblright{}.) Say it once more.
\end{tcolorbox}
